\section{Discussion}
\label{sec:discussion}

\paragraph{What the reconstruction result establishes.}
Correct Q/K-derived addresses improve held-out local response reconstruction
beyond global, token/channel, candidate/background, and matched shuffled
controls. The shuffled failures are important: they preserve table capacity
but break address--prototype semantics. E5 further makes unsupported addresses
visible through explicit fallback levels, while E7 shows that the useful signal
can be composed from smaller semantic subtables rather than a monolithic
Cartesian table.

\paragraph{Why the entry proxy is limited.}
The entry-count and FP32 value-array analyses isolate the storage associated
with supported scalar prototypes. They do not include keys, indices, counters,
padding, allocator overhead, memory traffic, or control logic. Consequently,
the observed 8.83\% and 20.54\% entry fractions support compact table
composition, not measured SRAM, latency, area, or energy savings. Those claims
require an implementation-level hardware or simulator study.

\paragraph{What current replacement adds.}
The end-to-end classifier results close part of the
reconstruction-to-replacement gap: calibrated Q/K lookup preserves QKFormer
accuracy across CIFAR-10/100 and $T=1/4$, and remains close to clean inference
when two early attention currents are replaced together. The aligned control
also remains strongest after calibration. The grouped binary-pattern
decomposition extends the result to all four attention projection currents,
but it uses more table values than the original projection weights and, by
itself, still targets projection currents rather than the whole runtime graph.
The later native audit closes the learned affine/BN/LIF-transition portion of
that gap, while residual additions, pooling, SSA matrix products/gating, and
global reductions remain outside the current replacement set. The evidence
therefore establishes strong QKFormer replacement fidelity for specified
operator classes, not measured hardware execution or unrestricted
cross-architecture transfer.

\paragraph{What the Spikformer mismatch probe establishes.}
The unmodified Spikformer intervention is retained as a compact compatibility
probe, not as the center of the transfer story. Learned temporal+channel
calibration recovers $9.88$ points over static lookup
($66.47\%\rightarrow76.35\%$), showing that calibration is useful even when the
SSA current geometry differs from QKFormer. The remaining gap identifies the
interaction contract that the positive QK-contract ablation then tests.
Missing unmodified-Spikformer semantic, robustness, and multi-current controls
remain visible as \tbd{} cells rather than being inferred from QKFormer.

\paragraph{What the QK-contract ablation adds.}
Replacing every SSA interaction with a QKFormer-like
$\operatorname{gate}(Q)\odot K$ contract preserves $87.54\%$ clean accuracy
and raises the learned temporal+channel lookup result to $86.16\%$, reducing
the original replacement drop by $10.43$ points. This is strong evidence that
lookup compatibility depends on interaction structure. It is not evidence
that the fine semantic address dominates all controls: shuffled address reaches
$85.83\%$ and token/channel mean reaches $86.18\%$ under the same
$3\times10^{-3}$ gate learning rate. Aligned is sensitive to correct address
assignment but remains indistinguishable from coarse token/channel state at
this accuracy resolution.

\paragraph{Scope of cross-architecture generalization.}
The QKFormer and QK-contract Spikformer results support bounded
cross-architecture generalization: TCSLU transfers across distinct backbones
when their attention interactions expose a compatible discrete QK contract.
The unmodified Spikformer probe shows that this capability is
architecture-conditioned. The newly opened CIFAR10-DVS run tests whether the
same contract remains useful on event data; evidence from another unmodified
attention family, together with matched controls and robustness checks, would
be required to broaden the claim beyond QK-compatible spiking attention.

\paragraph{From projection-current to native operator replacement.}
A follow-up first replaced all ten Q/K/V projections simultaneously with
scalar-level LUT decompositions. The local projection error was small
(aggregate NRMSE $1.20\times10^{-4}$), clipping was zero, and full-validation
Acc@1 fell by $0.45$ points. We therefore count the Q/K/V row as passing the
final project-level near-lossless criterion. MLP, patch-embedding, classifier,
and a cumulative substitution of all 32 convolutional/linear operators also
remained within gate.

The final audit moves beyond output hooks. It replaces all 32
Conv/Linear affine operators with native scalar lookup-and-accumulate modules,
folds 31 paired inference BN modules into those tables, and replaces all 35
LIF state transitions with finite transition LUTs. On CIFAR-100 full
validation, the resulting graph changes Acc@1 from $77.58\%$ to $77.38\%$ at
$T=1$ and from $81.02\%$ to $80.80\%$ at $T=4$, with aggregate LIF clipping
below $3\times10^{-9}$. This supports a strong completeness claim for the
learned affine, BN, and LIF-transition operator classes in the evaluated
QKFormer. It is not yet a complete full-graph QKFormer replacement: pooling,
residual addition, SSA matrix products/gating, and global/spatial/temporal
reductions remain ordinary arithmetic, and the table sizes remain fidelity
evidence rather than storage, latency, energy, or hardware-acceleration
measurements.
